\section{The institutions the ratio presumes}
\label{sec:7.2}
Here is what a pipeline would look like with the adverse party built in.

What stays unified is command, operational secrecy and the capacity to act at tempo. What is held adversely is the evidence bearing on civilian harm, the reviewer denominator, and the record. The evidence is composed by a body outside the chain of command, with its own sources, including the protected-site coordinates humanitarian organizations already submit for deconfliction, received with standing. The denominator is staffed and funded by that body, so it can't be cut the way the Central Command team was. The 2022 civilian-harm action plan was a first attempt at the separation; its dismantling is an argument for making it harder to dismantle, not evidence that it can't exist.

Models are formed by a public body that doesn't operate them, in the tradition of the public arsenals that made American arms before the contractor era, and licensed into the military on terms that include adverse review and a written refusal right. A line one vendor held by contract is then a standard term. In the stronger version, the same body pretrains the base models the schools form.

Precedents show authority outside the executive working inside real wars. Coalition “red card” holders in NATO operations could veto targets involving their assets, and Germany's Bundestag must approve each armed deployment. Both have limits I have to carry: a red card protects the vetoing state's interests and assets, which may not include civilians, and parliamentary approval is given per mission and never reaches a target. Both are nonetheless candidates for the adverse party the review ratio presumes, and a red-card holder is the rare candidate positioned to act at tempo. What a machine bearer adds here is least: acts assembled from permitted requests, cases no person sees, and the moment mid-act when an interrupt is still possible. It depends on the institutions for evidence and not for continuity: the deployment condition is met here by the same bodies whose presence narrows its role. The institutions it is compared with are people in a room composed so they can refuse.

Two offices are added for a reconstructed military: a civilian-harm body, which composes the evidence, demands proofs on logged outputs and receives the pre-registered ratio, and a model-forming public body, which licenses models on terms that include adverse review and a refusal right. The state appoints and funds both, only the floor court can remove either for cause, never the command the first reviews, and a licence the second has issued stands until the floor court rules. Appendix~\ref{sec:A} sets them out with the book's other offices.

\emph{What this doesn't do.} A reconstructed military accepts a speed disadvantage against rivals that don't reconstruct, and the floor costs a just defender more than that: its claims about incoming attacks are discounted, the action threshold sits above one, and the person-directed terms bar practices the law permits. What the defender buys with the price is a pipeline that will refuse the next government's war of choice as it refused the last one's, and how much speed that costs is left open.
